% Pursuing Stacks (cotes 134-2 to 134-8): archive cote 134-7, site page 5 (PDF page 6), typescript page 447.
% typescript. Draft compared: ps6.tex:185-244.
% AI-transcribed by Claude Opus 5.5 (claude-opus-5-5, adaptive thinking, effort high), Message Batches API, batch ps3 (msgbatch_01XdFmuVVBAcgc7EFxdtCqsj), 2026-11-12; not yet checked by eye.
\scanpage{134-7/5}
\mspage{447}

of affine schemes over k , thus we'll be working in the category of
functors (or presheaves over $\mathrm{Aff}_{/k}$ )
\[ \text{(\over{\unsure{g}}{f}')} \qquad (\mathrm{Aff}_{/k})^{o} \longrightarrow (\mathrm{Sets}) \quad , \]
more specifically U(k) will be a full subcategory of this category of
functors. We'll endow $\mathrm{Aff}_{/k}$ with one of the standard site structures,
the most convenient one here is the fpqc topology (faithfully flat qua-%
si compact topology), and work in the category of \uline{sheaves} in the
latter. In terms of the interpretation \over{\unsure{g}}{f}) as covariant functors\inserted{$X$} on
$\mathrm{Alg}_{/k}$ , this just means that we are restricting to functors which
1) commute to finite products and 2) are``compatible with faithfully
flat descent'' i.e. for any map
\[ k' \longrightarrow k'' \]
in $\mathrm{Alg}_{/k}$ such that k'' becomes a faithfully flat algebra over k' ,
the following diagram in (Sets)
\[ X(k') \longrightarrow X(k'') \rightrightarrows X(k''\otimes_{\unsure{k'}}k'') \]
is exact. Thus U(k) will be defined as a full subcategory of the\xout{x}
category of such functors, or ``sheaves''.

One way for defining U(k), is to \unsure{ppesent|present} it as the union of a
sequence of subcategories $U_m(k)$ (m a natural integer). We'll take
$U_o(k)$ to be just reduced to the final functors (i.e. X(k') is a one
point set for any k' in $\mathrm{Alg}_{/k}$ ), and define inductively $U_{m+1}(k)$ in
terms of $U_m(k)$ as follows. For any k-module M, let W(M) be the corres-%
ponding ``vector bundle'' , defined by
\[ \text{\hand{(g)}} \qquad W(M)(k') = M\otimes_k k' \quad , \]
then an object is in $U_{m+1}(k)$ iff there exists an object $X_m$ in $U_m(k)$,
and a k-module $M_{\xout{\unsure{m}}} = M_{m+1}$ , in such a way that $X=X_{m+1}$ should be iso-%
morphic to a ``torsor'' over $X_{m\xout{\unsure{+1}}}$ , with group W(M). I am not too sure
here whether we \unsure{shoubd|should} view the objects of $U_m(k)$ as endowed with the
extra structure consisting in giving the modules $M_1$, ... , $M_m$ used in
the inductive construction, and mo\unsure{re}over the successive fibrations -
if so, then of course the categories $U_m(k)$, and their union U(k),
will no longer be interpreted as a mere subcategory of the category of

\margin{5}
